\documentclass[11pt]{article}
\usepackage[left=0.85in,right=0.85in,top=0.65in,bottom=0.65in]{geometry}
\usepackage[T1]{fontenc}
\usepackage{lmodern}
\usepackage{amsmath,amssymb,amsthm,mathtools,booktabs}
\usepackage[colorlinks=true,linkcolor=blue,citecolor=blue,urlcolor=blue]{hyperref}
\usepackage{microtype}
\newtheorem{proposition}{Proposition}
\newcommand{\Q}{\mathbb Q}
\newcommand{\Z}{\mathbb Z}
\newcommand{\dd}{\,\mathrm{d}}
\newcommand{\HH}[2]{H_{#1}^{(#2)}}
\title{Explicit symmetric Pad\'e kernels\\and exact odd-zeta cancellation}
\author{}
\date{27 October 2026}
\begin{document}
\maketitle
\begin{abstract}
We identify the differential equations behind the moment data and construct
explicit symmetric rational kernels with Barnes-integral representations.
Finite rational linear conditions cancel every unwanted zeta coefficient.
An exact term ratio and tail bound yield a saddle-point limit for
normalized positive sums.
Two kernel families and a finite search over subsidiary conditions are
computed with complete gcd reduction. The construction produces pure-target
forms, including several small finite forms for $\zeta(5)$, but does not
establish a sequence of nonzero primitive forms tending to zero.
The all-odd irrationality theorem is not proved.
\end{abstract}

\section{Which differential equation is actually available?}
The existing coefficients are defined by
\[
 \Delta_{N,a}(X)=\det[X-\HH{a+i+j}{s}]_{0\le i,j<N}=A_{N,a}X+B_{N,a},
 \quad Q=\operatorname{lcm}(1,\ldots,a+2N-2)^s,
\]
\[
 U=Q^NA_{N,a},\qquad V=Q^NB_{N,a},\qquad
 g=\gcd(|U|,|V|),\qquad u=U/g.
\]
An empty lcm is one. These formulas define $u$ arithmetically; no linear
differential equation or single Barnes integral for the gcd quotient
has been established.

There is an exact equation for the moment-generating data. For $|z|<1$,
absolute convergence gives
\[
 \sum_{n\ge0}\HH{n}{s}z^n=\frac{\operatorname{Li}_s(z)}{1-z},
 \qquad \operatorname{Li}_s(z)=\sum_{n\ge1}\frac{z^n}{n^s}.
\]
With $\theta=z\,d/dz$, we have
$\theta^s\operatorname{Li}_s=z/(1-z)$, and hence
\begin{equation}
 \bigl((1-z)\theta-1\bigr)\theta^s\operatorname{Li}_s=0.
 \label{eq:polylog}
\end{equation}
For $s\ge2$, its indicial exponents are
\[
 \begin{array}{c|c}
 0&0\ (s\text{ times}),\ 1\\
 1&0,1,\ldots,s-2,\ s-1,\ s-1\\
 \infty&0\ (s+1\text{ times})
 \end{array}
\]
Indeed, the indicial polynomials are respectively
$\rho^s(\rho-1)$, a nonzero multiple of
$\rho(\rho-1)\cdots(\rho-s+1)(\rho-s+1)$, and a nonzero
multiple of $\rho^{s+1}$ after setting $z=1/w$.
These exponents do not determine an integral period lattice.

\section{The exact arithmetic price of subtracting a counterterm}
Let $\operatorname{den}(q)>0$ denote a reduced rational denominator.
For rational $x,c$, elementary denominator arithmetic gives
\begin{equation}
 \operatorname{den}(x-c)\mid\operatorname{den}(x)\operatorname{den}(c),
 \qquad
 \operatorname{den}(x)\mid\operatorname{den}(x-c)\operatorname{den}(c).
 \label{eq:center}
\end{equation}
The first follows by subtraction; the second by $x=(x-c)+c$.
Consequently, subtracting a rational counterterm with denominator $m$
can reduce the denominator of $x$ by at most a factor $m$.
An integer counterterm leaves it unchanged. Lean verifies these statements.

Euler--Maclaurin expansions can sharply estimate analytic remainders.
To use a subtraction in an integer linear form, its rational denominator
must still enter \eqref{eq:center}. A nonrational logarithmic counterterm
adds another period unless its coefficient is independently canceled.
Merely introducing and subtracting an intermediate scale does not alter
the original reduced coefficient.

\section{A concrete Barnes family with subsidiary conditions}
Fix $r,n\ge1$, put $s=2r+1$, $A=2r$, and choose $0\le m\le rn$.
These kernels have a different arithmetic budget from the single-shift
determinants; the earlier lower bound is not asserted for them. Define
\begin{equation}
 R_{r,n,m}(t)=
 \frac{\prod_{v=1}^{m}(t-v)(t+n+v)}
      {\prod_{j=0}^{n}(t+j)^A}.
 \label{eq:kernel}
\end{equation}
This is a proper rational function of decay at least $t^{-2r}$.
Away from the apparent gamma poles its exact gamma-ratio representation is
\[
 R_{r,n,m}(t)=
 \frac{\Gamma(t)^{A+1}\Gamma(t+n+m+1)}
      {\Gamma(t-m)\Gamma(t+n+1)^{A+1}}.
\]
Both sides continue as the same rational function. Direct cancellation gives
\begin{align}
 R_{r,n,m+1}(t)&=
 [t(t+n)-(m+1)(n+m+1)]R_{r,n,m}(t),\\
 \frac{R_{r,n,m}(t+1)}{R_{r,n,m}(t)}
 &=\frac{t^{A+1}(t+n+m+1)}
         {(t-m)(t+n+1)^{A+1}},
 \label{eq:ratio}
\end{align}
where the second identity is used away from zeros and poles.
These are actual recurrences for the rational kernel.

\begin{proposition}
For $0<c<1$,
\begin{equation}
 \Lambda_{r,n,m}:=-\sum_{k=1}^{\infty}R'_{r,n,m}(k)
 =\frac{1}{2\pi i}\int_{c-i\infty}^{c+i\infty}
       \left(\frac{\pi}{\sin\pi t}\right)^2R_{r,n,m}(t)\dd t.
 \label{eq:barnes}
\end{equation}
Moreover $\Lambda_{r,n,m}\in
\Q+\Q\zeta(3)+\cdots+\Q\zeta(2r+1)$, with explicitly computable
rational coefficients.
\end{proposition}
\begin{proof}
The series converges absolutely by the decay of $R'$.
The integrand in \eqref{eq:barnes} has residue $R'(k)$ at each positive
integer. Close the upward vertical line to the right along half-integer
rectangles. The orientation is clockwise; the remaining edges vanish by
rational decay and the exponential decay of $\sin(\pi t)^{-2}$.
This proves the sign and normalization of \eqref{eq:barnes}.

Write $R=\sum_{j=0}^n\sum_{k=1}^A c_{j,k}(t+j)^{-k}$.
Its reflection symmetry $R(-n-t)=R(t)$ gives
$c_{n-j,k}=(-1)^kc_{j,k}$. Differentiating the finite expansion and using
$\sum_{\ell\ge1}(\ell+j)^{-q}=\zeta(q)-H_j^{(q)}$ gives
\[
 \Lambda=\sum_{q=2}^{A+1}(q-1)\left(\sum_jc_{j,q-1}\right)\zeta(q)
          -\sum_{j=0}^n\sum_{k=1}^A k c_{j,k}\HH{j}{k+1}.
\]
Even $q$ vanish by the reflection identity.
The coefficients follow from the finite Taylor formula
\[
 c_{j,k}=\frac1{(A-k)!}
 \left.\frac{d^{A-k}}{dt^{A-k}}[(t+j)^A R(t)]\right|_{t=-j}.
\]
\end{proof}

Take $r$ block lengths $m_1,\ldots,m_r$, and denote their coefficients of
$\zeta(2q+1)$ by $a_{q,j}$. Choose rational, then integer, weights satisfying
\[
 \sum_{j=1}^r w_j a_{q,j}=0\qquad(1\le q<r).
\]
There is a nonzero rational solution because there are $r-1$ equations.
The target coefficient must separately be nonzero; the finite audits check
both rank and this nonvanishing. The resulting form is
\[
 L=\sum_jw_j\Lambda_{r,n,m_j}=a\zeta(s)+b.
\]
For consecutive block lengths, the combined kernel is
$R_{r,n,m_{\min}}(t)P(t(t+n))$ with $\deg P\le r-1$.
Thus the subsidiary conditions define a concrete polynomial deformation.

A second audited family replaces \eqref{eq:kernel} by
\[
 S_{r,n,m}(t)=
 \frac{(2t+n)^\epsilon\prod_{v=1}^{m}(t-v)(t+n+v)}
      {\prod_{j=0}^{n}(t+j)^s},\qquad \epsilon\equiv n\pmod2,
\]
with $2m+\epsilon\le s(n+1)-2$.
Then $S(-n-t)=-S(t)$, so $\sum_{k\ge1}S(k)$ again has only odd zeta
terms and a rational constant. Its total simple-pole residue is zero
by proper decay; this accounts for the convergent first-order pole terms.

There is also an analytic nonvanishing mechanism. At positive integers,
$S(k)=0$ for $k\le m$ and $S(k)>0$ for $k>m$: beyond the numerator
zeros every displayed factor is positive. Lean verifies these inequalities
for the actual rational kernel and the positive-value conclusion for any
of its convergent sums. If the subsidiary weights have one sign and are
not all zero, the combined integer summands have that sign and are
strictly nonzero beyond all block lengths. Their sum is therefore nonzero.
Constructing such weights with adequate size bounds remains an arithmetic
requirement; positivity is not assumed for arbitrary nullvectors.

\subsection{The functional Pad\'e bridge}
Let $c_{j,q}$ be the partial-fraction coefficients of $S$. Define
\[
 P_q(z)=\sum_{j=0}^n c_{j,q}z^{n-j},\qquad
 P_0(z)=\sum_{j=0}^n\sum_{q=1}^s c_{j,q}
              \sum_{\ell=1}^j\frac{z^{n-j+\ell}}{\ell^q}.
\]
All these polynomials have degree at most $n$. Reindexing absolutely
convergent series for $|z|<1$ gives the exact functional identity
\begin{equation}
 \sum_{q=1}^sP_q(z)\operatorname{Li}_q(z)-P_0(z)
 =z^n\sum_{k\ge1}S(k)z^k=O(z^{n+m+1}).
 \label{eq:functional}
\end{equation}
The zero order follows from the first $m$ numerator zeros.
Reflection gives $P_q(1)=0$ for even $q$, and proper decay gives
$P_1(1)=0$. Hence the logarithmic term
$P_1(z)\operatorname{Li}_1(z)$ tends to zero at $z=1$:
$P_1(z)=O(1-z)$ and $\operatorname{Li}_1(z)=-\log(1-z)$.
The additional odd-coefficient conditions cancel the remaining unwanted
values. This supplies an actual functional Pad\'e problem at $0$ and
subsidiary conditions at $1$, rather than assuming an identification with
the polylogarithm differential system.

\section{What the cancellations achieve, and their remaining cost}
At $r=2,n=1$, choosing the two weights $-23,-145$ for block lengths
$2,1$ gives, with $T=t(t+1)$,
\[
 K(t)=\frac{-23T^2+39T+14}{T^4}
 =-\frac{17}{t^3}+\frac{14}{t^4}
   +\frac{17}{(t+1)^3}+\frac{14}{(t+1)^4}.
\]
Consequently
\[
 -K'(t)=-\frac{51}{t^4}+\frac{56}{t^5}
         +\frac{51}{(t+1)^4}+\frac{56}{(t+1)^5},
 \qquad -\sum_{k\ge1}K'(k)=112\zeta(5)-107.
\]
Lean verifies the rational identities and the corresponding
\texttt{HasSum} identity for the actual real zeta value.
It also verifies $112\zeta(5)-107\ge5$.
Thus exact lower-period cancellation alone does not give smallness.

The two families were checked at $92$ weight/index/profile choices.
A separate search enumerated $2140$ distinct nonzero-target two-kernel
forms for $\zeta(5)$, with complete gcd reduction.
The indices are $2,4,8,12,20$, with respectively
$21,66,231,496,1326$ distinct candidates.
For the first profile the full combined kernel is explicitly
\[
 F(t)=601088S_{2,2,2}(t)+37632S_{2,2,3}(t)
 =\frac{256(147t^2+294t+143)(t-1)(t-2)(t+3)(t+4)}
 {t^5(t+1)^5(t+2)^5}.
\]
Its partial fractions give
$\sum_{k\ge1}F(k)=91776\zeta(5)-95165>0$, with the entire
$\zeta(3)$ coefficient canceled. The next best primitive form is
$1216512\zeta(5)-1261435$.
Both absolute values are certified below one. All signs and finite
comparisons use exact rational zeta enclosures.
Finite smallness is not a limit theorem.

If $D$ clears the denominators of $a,b$ and
$g=\gcd(|Da|,|Db|)$, the necessary combined estimate is
\[
 0<|L_{\rm prim}|=\frac Dg|L|
 \le\frac Dg\sum_j|w_j|\,|\Lambda_j|\longrightarrow0.
\]
A Barnes saddle estimate must therefore be paired with bounds on the
subsidiary-condition weights and the exact arithmetic normalization.
These bounds have not been proved to win for every odd target.

The general contour, gamma-ratio, indicial, functional Pad\'e, and
saddle-limit arguments are paper proofs here. Lean verifies the centering budget, finite period-coordinate
elimination, the explicit $\zeta(5)$ series identities, and integer-kernel
positivity. No all-index target nonvanishing or winning arithmetic estimate
is asserted.

There are specialized alternatives to a general Siegel--Shidlovskii
statement for $G$-functions. Fischler's multipoint multiplicity argument
\cite{Fischler} proves the infinitely-many-odd-values result, whereas
Rivoal \cite{Rivoal} shows that the direct general $G$-function analogue
is false. Recent Hurwitz-zeta Pad\'e asymptotics \cite{Kandaian} assume
determinantal nondegeneracy and a regular equilibrium regime; they do
not assert the integer coefficient estimate above.
The explicit families here provide objects on which to test such methods,
without assuming their missing arithmetic conclusion.

\section{Analytic control of the positive summands}
This calculation concerns the sums before the subsidiary-condition
combination and final primitive normalization. Restrict to even $n$,
fix $0<\beta<s/2$, and let the integers $m=m_n$ satisfy $m_n/n\to\beta$. Put
\[
 F_{n,m}(k)=\frac{(n!)^s}{(m!)^2}S_{r,n,m}(k),\qquad
 {\cal A}_{n,m}=\sum_{k>m}F_{n,m}(k).
\]
The summands are positive and have the exact ratio
\[
 \frac{F_{n,m}(k+1)}{F_{n,m}(k)}
 =\frac{k^{s+1}(k+n+m+1)}{(k-m)(k+n+1)^{s+1}}.
\]
Define, for $x>\beta$,
\[
 \begin{aligned}
 \Phi_{s,\beta}(x)={}&(s+1)[x\log x-(x+1)\log(x+1)]\\
 &+(x+1+\beta)\log(x+1+\beta)\\
 &-(x-\beta)\log(x-\beta)-2\beta\log\beta.
 \end{aligned}
\]
There is a unique maximizing point $x_\beta$. Indeed, with
$v=x/(x+1)$, the equation $\Phi'=0$ becomes
\[
 G(v):=\beta(1+v^{s+1})-\sum_{j=1}^s v^j=0.
\]
For $0<v<1$,
$G'(v)=\beta(s+1)v^s-\sum_{j=1}^sjv^{j-1}<0$, because
$\sum jv^{j-1}\ge s(s+1)v^s/2$ and $2\beta<s$.
Also $G(\beta/(1+\beta))>0$ and $G(1)<0$.
The same argument, with $\beta$ replaced by $m/(n+1)$, proves
unimodality of the exact discrete summands.

Here is a uniform tail bound, so the infinite sum introduces no missing
exponential factor. Set
\[
 E=s(n+1)-2m-2,\quad B=2m^2+sn(n+1),\quad
 K=\max\{2m,m+1,\lceil B/E\rceil\}.
\]
For large $n$, $E>0$ and $K=O(n)$. For real $t\ge K$,
\[
 \frac{d}{dt}\log[t^2S(t)]
 \le-\frac E t+\frac B{t^2}\le0.
\]
To verify this, use
$\sum_{v=1}^m(t-v)^{-1}\le m/t+2m^2/t^2$,
$\sum_{v=1}^m(t+n+v)^{-1}\le m/t$, and
$\sum_{j=0}^n(t+j)^{-1}\ge(n+1)/t-n(n+1)/t^2$.
Thus $F(k)\le F(K)(K/k)^2$ for $k\ge K$, and
\[
 \max_{k>m}F(k)\le {\cal A}_{n,m}
 \le(2K+1)\max_{k>m}F(k).
\]
Finally the elementary factorial estimate
$\log(q!)=q\log q-q+O(\log(q+2))$, obtained by bounding
the sum of $\log j$ between its adjacent integrals, gives
$n^{-1}\log F_{n,m}(\lfloor nx\rfloor)\to\Phi_{s,\beta}(x)$
locally uniformly for $x>\beta$. The exact ratio locates the peak at
$k/n\to x_\beta$. Consequently
\begin{equation}
 \lim_{\substack{n\to\infty\\ n\ \mathrm{even}}}
 \frac1n\log{\cal A}_{n,m_n}
 =\Phi_{s,\beta}(x_\beta).
 \label{eq:analyticrate}
\end{equation}
In particular the certified rate for $s=5$, $\beta=1$ lies in
$(-4.970794,-4.970793)$. This proves exponential decay of those normalized
positive sums. The scripts certify $44$ negative saddle-rate intervals
and $22$ finite positive-sum enclosures against independent period
enclosures.

For $\beta=1$, decay follows at every odd weight, not just the tested
weights. For $x>1$, $B(x):=x\log x-(x+1)\log(x+1)<0$, and hence
\[
 \Phi_{s,1}(x)=\Phi_{3,1}(x)+(s-3)B(x)\le\Phi_{3,1}(x).
\]
The certified base maximum is less than $-1.99$. Thus, for each fixed
odd $s\ge3$, $(n!)^{s-2}\sum_{k\ge1}S_{r,n,n}(k)\le e^{-1.99n}$
for all sufficiently large even $n$.

\subsection{Completion of the lower-period analytic step}
For a fixed $s=2r+1$, use the $r$ normalized sums with block lengths
$m_j=n-j$, $0\le j<r$. Each has rate limit $\Phi_{s,1}(x_1)$, so
all are bounded by $e^{-1.99n}$ for sufficiently large even $n$.
Write their lower-period coefficient matrix as $C_n\in
M_{r-1,r}(\Q)$. Dimension gives a nonzero $w_n\in\ker C_n$.
These weights are computable: when $C_n$ has full row rank, use the
signed maximal minors $w_{n,j}=(-1)^j\det C_n^{\widehat j}$.
Laplace expansion with a repeated row proves $C_nw_n=0$.
Otherwise rational row reduction and a free coordinate produce a kernel vector.
Put $M_n=\sum_j|w_{n,j}|>0$ and $\widehat w_{n,j}=w_{n,j}/M_n$.
Then, exactly,
\[
 C_n\widehat w_n=0,\qquad
 \sum_j|\widehat w_{n,j}|=1,\qquad
 L_n:=\sum_j\widehat w_{n,j}{\cal A}_{n,n-j}
      =a_n\zeta(s)+b_n,
\]
\[
 |L_n|\le\sum_j|\widehat w_{n,j}|\,|{\cal A}_{n,n-j}|
       \le e^{-1.99n}.
\]
This completes lower-period cancellation and preserves analytic decay.
Lean verifies the rational kernel-vector existence, normalization,
cancellation, norm bound, and resulting convergence implication.

This normalization does not itself give small integer forms.
If $a_n\ne0$, multiplying both rational coefficients by the same nonzero
rational does not change the reduced denominator of $b_n/a_n$; Lean verifies
this invariance. One must still prove nonvanishing and
$(D_n/g_n)|L_n|\to0$ after the full integer normalization.

\begin{thebibliography}{9}
\small
\bibitem{Rivoal}
T. Rivoal, \emph{Remarks on the impossibility of a Siegel--Shidlovskii
like theorem for G-functions}, Hardy--Ramanujan Journal 38 (2015).
\url{https://hrj.episciences.org/1357}.
\bibitem{Fischler}
S. Fischler, \emph{Shidlovsky's multiplicity estimate and Irrationality
of zeta values}, 2016.
\url{https://arxiv.org/abs/1609.09770}.
\bibitem{Kandaian}
A. Kandaian, \emph{Multipoint Pad\'e Approximation of the Hurwitz Zeta
Function and a Riemann--Hilbert Steepest Descent Analysis}, preprint, 2026.
\url{https://arxiv.org/abs/2602.08399}.
\end{thebibliography}
\end{document}
